% ------------------------------------------------------------------------
% file `3_thales_similitudes-exercise-3-solution-body.tex'
%   in folder `xsim/'
%
%     solution of type `exercise' with id `3'
%
% generated by the `solution' environment of the
%   `xsim' package v0.21 (2022/02/12)
% from source `3_thales_similitudes' on 2023/06/07 on line 127
% ------------------------------------------------------------------------
    \begin{tasks}
        \task
        $
            \begin{aligned}[t]
                CD \sslash AB \Rightarrow
                \dfrac{\abs{EC}}{\abs{EA}} & = \dfrac{\abs{CD}}{\abs{AB}}
                \qquad\text{Par le théorème de Thalès}                    \\
                \dfrac{1,5}{4}             & = \dfrac{x}{3}               \\
                x                          & = 3\cdot\dfrac{1,5}{4}       \\
                x                          & = \dfrac{4,5}{4}             \\
                x                          & = 1,125
            \end{aligned}
        $
        \task $ABC \sim DEF$ ?
        $
            \begin{aligned}[t]
                \dfrac{\abs{AC}}{\abs{DF}}
                 & = \dfrac{2}{4,2} \\
                 & = \dfrac{20}{42} \\
                 & = \dfrac{10}{21} \\
            \end{aligned}
            \hfil
            \begin{aligned}[t]
                \dfrac{\abs{AB}}{\abs{DE}}
                 & = \dfrac{3}{6,3} \\
                 & = \dfrac{30}{63} \\
                 & = \dfrac{10}{21} \\
            \end{aligned}
            \hfil
            \begin{aligned}[t]
                \dfrac{\abs{BC}}{\abs{EF}}
                 & = \dfrac{4}{8,4} \\
                 & = \dfrac{40}{84} \\
                 & = \dfrac{10}{21}
            \end{aligned}
        $\bigbreak

        Par le cas CCC, $ABC \sim DEF$. Donc $\abs{\widehat{D}} = \abs{\widehat{A}} = \ang{104}$.
        \task
        $
            \begin{aligned}[t]
                CD \sslash AB \Rightarrow
                \dfrac{\abs{AE}}{\abs{ED}} & = \dfrac{\abs{EB}}{\abs{EC}}
                \qquad\text{Par le théorème de Thalès}                    \\
                \dfrac{x}{9}               & = \dfrac{3}{5}               \\
                x                          & = 9\cdot\dfrac{3}{5}         \\
                x                          & = \dfrac{18}{5}              \\
                x                          & = 3,6
            \end{aligned}
        $

        \task \(ABC \sim DEF \) ?
        \begin{itemize}
            \item \(\dfrac{\abs{EF}}{\abs{AC}}=\dfrac{\abs{DF}}{\abs{BC}}=\dfrac{3}{5}\)
            \item \(ABC\) est un triangle isocèle en \(C\) et \(\abs{\widehat{A}}=\ang{33,5}\) donc,
                  \begin{align*}
                      \abs{\widehat{C}} & =\ang{180}-2\cdot \ang{33,5} \\
                                        & =\ang{113}                   \\
                                        & =\abs{\widehat{F}}
                  \end{align*}
        \end{itemize}
        Par le cas CAC, \(ABC \sim DEF\).

        On en déduit,
        \begin{align*}
            \dfrac{x}{5} & = \dfrac{1,8}{3}         \\
            x            & = \dfrac{1,8}{3} \cdot 5 \\
            x            & = 3
        \end{align*}
    \end{tasks}
